\documentclass[11pt,a4paper]{article}

\usepackage[T1]{fontenc}
\usepackage[utf8]{inputenc}
\usepackage{lmodern}
\usepackage[margin=1in]{geometry}
\usepackage{amsmath,amssymb,amsthm}
\usepackage{booktabs}
\usepackage[hidelinks]{hyperref}
\usepackage{bookmark}

\setlength{\emergencystretch}{1.5em}
\renewcommand{\arraystretch}{1.18}

\newtheorem{theorem}{Theorem}[section]
\newtheorem{proposition}[theorem]{Proposition}
\newtheorem{lemma}[theorem]{Lemma}
\newtheorem{corollary}[theorem]{Corollary}
\theoremstyle{definition}
\newtheorem{definition}[theorem]{Definition}
\theoremstyle{remark}
\newtheorem{remark}[theorem]{Remark}

\newcommand{\Cl}{\mathrm{Cl}}
\newcommand{\Ann}{\mathrm{Ann}}
\newcommand{\codim}{\mathrm{codim}}
\newcommand{\rank}{\mathrm{rank}}
\newcommand{\rad}{\mathrm{rad}}
\newcommand{\GL}{\mathrm{GL}}

\title{\textbf{Disproof of the Zero-Divisor Metric Rank-Deficit Law\\
(Conjecture 00000007621)}}
\author{Formalized in Core Lean 4.19.0}
\date{November 2026}

\begin{document}
\maketitle

\begin{abstract}
We disprove Conjecture~00000007621 (\emph{zero-divisor metric rank-deficit law}), which asserts that the zero-divisor variety defined by multiplicative annihilation obeys an explicit dimension law whose governing constant is the metric rank deficit (equal to the metric nullity count). While the linear-algebraic identification between metric rank deficit $n - \rank(B_Q)$ and metric nullity $r = \dim \ker(B_Q)$ is tautological, the dimension of the zero-divisor variety $Z(A)$ of a metric algebra $A = \Cl(V, Q)$ is neither determined by $(n, r)$ nor governed by $r$. Specifically, we prove a complete radical-quotient and Bott-periodicity theorem showing that $\dim_{\mathbb{R}} Z(\Cl_{p,q,r}(\mathbb{R})) = 2^n - c(p - q \bmod 8)$ with $c \in \{1, 2, 4\}$ equal to the real dimension of the underlying division ring $\mathbb{D} \in \{\mathbb{R}, \mathbb{C}, \mathbb{H}\}$, completely independent of the metric nullity $r$, while over any algebraically closed field $K$ of characteristic $\neq 2$, $\dim_K Z(\Cl(K^n, Q)) = 2^n - 1$ for every $r \in \{0, \dots, n\}$. Already for nondegenerate metrics ($r = 0$), $\dim_{\mathbb{R}} Z$ takes distinct values $\{0, 1\}$ at $n = 1$, $\{0, 3\}$ at $n = 2$, and $\{4, 6, 7\}$ at $n = 3$. We further prove that restricting to grade-$1$ vectors $Z(A) \cap V$, passing to finite-field point counts over $\mathbb{F}_q$, passing to multiplicative annihilator dimensions $\dim \Ann_R(x)$, or passing to incidence pairs and Jordan spin factors likewise refutes the conjecture. All core algebraic classifications, radical invariance theorems, annihilator separations, and exact $\mathbb{F}_q$ point-count formulas are formally verified from first principles in Core Lean~4.19.0 with zero external dependencies and zero custom axioms.
\end{abstract}

\section{Problem Statement and Algebraic Framework}

Conjecture~00000007621 in \texttt{SOURCE.md} states (in parallel English and Chinese formulations):
\begin{quote}
\emph{Definition: Zero divisors as multiplicative annihilation structure. Conjecture: The dimension law of the zero-divisor variety is explicit; the constant of the law is the metric rank deficit; and the constant of the metric rank deficit is the nullity count. (zero-divisor metric rank-deficit law)}
\end{quote}

Let $K$ be a field and let $(V, Q)$ be an $n$-dimensional quadratic space over $K$ (with $\mathrm{char}\,K \neq 2$ when passing between quadratic and symmetric bilinear forms), where
\[
B_Q(u, v) = \frac{1}{2}\bigl(Q(u + v) - Q(u) - Q(v)\bigr)
\]
is the associated symmetric bilinear metric. Its radical is
\[
\rad(B_Q) = \ker(B_Q) = \{v \in V : B_Q(v, w) = 0 \text{ for all } w \in V\}.
\]
By the rank--nullity theorem for the linear map $V \to V^*$ induced by $B_Q$, the \emph{metric rank deficit} equals the \emph{metric nullity count}:
\begin{equation}\label{eq:rank-nullity}
n - \rank(B_Q) = \dim_K \ker(B_Q) =: r \in \{0, 1, \dots, n\}.
\end{equation}
Thus the final clause of the conjecture merely restates~\eqref{eq:rank-nullity}. The substantive content of Conjecture~00000007621 is the claim that in the multiplicative algebra attached to $(V, Q)$---canonically the Clifford algebra $A = \Cl(V, Q)$ generated by $V$ subject to $v^2 = Q(v)\cdot 1$ (including the exterior algebra $\Lambda(V) = \Cl(V, 0)$ at $r = n$)---the dimension law of the zero-divisor variety is governed by the constant $r$.

\begin{definition}[Multiplicative annihilation and zero-divisor locus]
Let $A$ be a finite-dimensional unital associative $K$-algebra. For $x \in A$, the right and left multiplicative annihilators are
\[
\Ann_R(x) = \ker(L_x) = \{y \in A : xy = 0\}, \qquad
\Ann_L(x) = \ker(R_x) = \{y \in A : yx = 0\},
\]
where $L_x(y) = xy$ and $R_x(y) = yx$. The \emph{zero-divisor variety} (or zero-divisor locus) of $A$ is
\begin{align*}
Z(A) &= \{x \in A : \Ann_R(x) \neq \{0\} \text{ or } \Ann_L(x) \neq \{0\}\} \\
     &= \{x \in A : \exists\, y \in A \setminus \{0\},\; xy = 0 \text{ or } yx = 0\}.
\end{align*}
\end{definition}

\begin{lemma}[Unit--zero-divisor dichotomy]\label{lem:dichotomy}
Let $A$ be a finite-dimensional unital associative $K$-algebra with $1 \neq 0$. For every $x \in A$, the following are equivalent:
\begin{enumerate}
\setlength{\itemsep}{2pt}\setlength{\parskip}{0pt}
\item[(i)] $x \in Z(A)$;
\item[(ii)] $\det(L_x) = 0$;
\item[(iii)] $x \notin A^\times$ (that is, $x$ has no two-sided multiplicative inverse in $A$).
\end{enumerate}
In particular, $Z(A) = A \setminus A^\times = \{x \in A : \det(L_x) = 0\}$ is a Zariski-closed affine cone in $A$.
\end{lemma}

\begin{proof}
If $x \in A^\times$ with $x x^{-1} = x^{-1} x = 1$, then $xy = 0 \implies y = x^{-1}(xy) = 0$ and $yx = 0 \implies y = (yx)x^{-1} = 0$, so $x \notin Z(A)$. Conversely, suppose $x \notin Z(A)$. Then $\ker(L_x) = \Ann_R(x) = \{0\}$, so $\det(L_x) \neq 0$. Because $\dim_K A < \infty$, the injective linear endomorphism $L_x : A \to A$ is surjective, yielding $y \in A$ with $xy = 1$. Then $yz = 0 \implies z = (xy)z = x(yz) = 0$, so $L_y$ is also injective and hence surjective, giving $w \in A$ with $yw = 1$. Associativity gives $x = x(yw) = (xy)w = w$, so $yx = 1$ and $x \in A^\times$. Finally, $\det(L_x) = 0 \iff \ker(L_x) \neq \{0\}$, which by the same argument is equivalent to $x \notin A^\times$.
\end{proof}

\section{Exact Dimension Law and Refutation on Clifford Algebras}

We first prove the exact structural theorem governing $Z(\Cl(V, Q))$ in the presence of arbitrary metric rank deficit $r \ge 0$, demonstrating that $r$ cancels completely from the codimension of $Z(\Cl(V, Q))$.

\begin{theorem}[Radical projection invariance]\label{thm:radical}
Let $(V, Q)$ be an $n$-dimensional quadratic space over a field $K$ with $\mathrm{char}\,K \neq 2$, let $R = \ker(B_Q)$ have dimension $r$, and choose an orthogonal decomposition $V = W \perp R$ so that $\bar{Q} = Q|_W$ is nondegenerate of rank $k = n - r$. Let $\pi : \Cl(V, Q) \to \Cl(W, \bar{Q})$ be the canonical surjective algebra homomorphism with kernel $J = \Cl(V, Q)\cdot R\cdot \Cl(V, Q)$. Then:
\begin{enumerate}
\setlength{\itemsep}{2pt}\setlength{\parskip}{0pt}
\item[(i)] $J^{r+1} = \{0\}$, and $\dim_K J = 2^n - 2^{n-r}$;
\item[(ii)] For every $x \in \Cl(V, Q)$, $x \in Z(\Cl(V, Q)) \iff \pi(x) \in Z(\Cl(W, \bar{Q}))$;
\item[(iii)] As an affine algebraic variety (over $K = \mathbb{R}$, $K = \mathbb{C}$, or any field $K$),
\[
Z\bigl(\Cl(V, Q)\bigr) = \pi^{-1}\!\Bigl(Z\bigl(\Cl(W, \bar{Q})\bigr)\Bigr)
\cong Z\bigl(\Cl(W, \bar{Q})\bigr) \times K^{2^n - 2^{n-r}},
\]
and hence $\codim_K Z(\Cl(V, Q)) = \codim_K Z(\Cl(W, \bar{Q}))$ is independent of $r$.
\end{enumerate}
\end{theorem}

\begin{proof}
Choose an orthogonal basis $w_1, \dots, w_k$ of $W$ and a basis $g_1, \dots, g_r$ of $R$. In $\Cl(V, Q)$, the generators $g_1, \dots, g_r$ satisfy $g_j^2 = Q(g_j) = 0$ and anticommute pairwise as well as with every $w_i$. Every element of $J$ is a $K$-linear combination of basis blades containing at least one factor $g_j$. Any product of $r + 1$ such blades can be rearranged up to sign so that at least $r + 1$ factors from $\{g_1, \dots, g_r\}$ stand adjacent; by the pigeonhole principle, some $g_j$ occurs at least twice and vanishes since $g_j^2 = 0$. Hence $J^{r+1} = \{0\}$, and counting basis blades with and without $g$-factors gives $\dim_K J = 2^n - 2^k$.

For~(ii), by Lemma~\ref{lem:dichotomy} it suffices to prove $x \in \Cl(V, Q)^\times \iff \pi(x) \in \Cl(W, \bar{Q})^\times$. The forward implication holds because $\pi$ is a unital algebra homomorphism. Conversely, if $\pi(x)$ has inverse $\bar{y} \in \Cl(W, \bar{Q})^\times$, choose $y \in \Cl(V, Q)$ with $\pi(y) = \bar{y}$. Then $\pi(1 - xy) = 0$, so $\eta := 1 - xy \in J$ satisfies $\eta^{r+1} = 0$. Therefore
\[
xy = 1 - \eta \in \Cl(V, Q)^\times \quad\text{with inverse}\quad u = \sum_{m=0}^r \eta^m,
\]
giving $x(yu) = 1$. By Lemma~\ref{lem:dichotomy}, a right inverse in a finite-dimensional unital associative algebra is automatically a two-sided inverse, so $x \in \Cl(V, Q)^\times$. Claim~(iii) follows immediately from linear splitting $\Cl(V, Q) \cong \Cl(W, \bar{Q}) \oplus J$.
\end{proof}

Over $\mathbb{R}$, write $\Cl_{p,q,r}(\mathbb{R})$ for the Clifford algebra of the quadratic form
\[
Q_{p,q,r}(x_1,\dots,x_p, y_1,\dots,y_q, z_1,\dots,z_r) = \sum_{i=1}^p x_i^2 - \sum_{j=1}^q y_j^2 + \sum_{\ell=1}^r 0\cdot z_\ell^2
\]
on $\mathbb{R}^n$, where $n = p + q + r$, $\rank(B_Q) = p + q$, and the metric rank deficit (nullity) is $r$.

\begin{theorem}[Exact real Bott-periodic dimension law]\label{thm:real-law}
For all integers $p, q, r \ge 0$ with $n = p + q + r$, the real algebraic dimension of $Z(\Cl_{p,q,r}(\mathbb{R}))$ is
\begin{equation}\label{eq:bott-law}
\dim_{\mathbb{R}} Z\bigl(\Cl_{p,q,r}(\mathbb{R})\bigr) = 2^n - c(p - q \bmod 8), \qquad
c(k) = \begin{cases}
1, & k \equiv 0, 1, 2 \pmod 8,\\
2, & k \equiv 3, 7 \pmod 8,\\
4, & k \equiv 4, 5, 6 \pmod 8.
\end{cases}
\end{equation}
\end{theorem}

\begin{proof}
By Theorem~\ref{thm:radical}(iii), $\codim_{\mathbb{R}} Z(\Cl_{p,q,r}(\mathbb{R})) = \codim_{\mathbb{R}} Z(\Cl_{p,q}(\mathbb{R}))$. By the Atiyah--Bott--Shapiro classification of nondegenerate real Clifford algebras under $v^2 = Q(v)\cdot 1$, $\Cl_{p,q}(\mathbb{R})$ is isomorphic to $M_N(\mathbb{D})$ when $p - q \not\equiv 1, 5 \pmod 8$ and to $M_N(\mathbb{D}) \oplus M_N(\mathbb{D})$ when $p - q \equiv 1, 5 \pmod 8$, where the underlying real division algebra $\mathbb{D}$ is:
\begin{itemize}
\setlength{\itemsep}{2pt}\setlength{\parskip}{0pt}
\item $\mathbb{D} = \mathbb{R}$ ($\dim_{\mathbb{R}} \mathbb{D} = 1$) for $p - q \equiv 0, 1, 2 \pmod 8$;
\item $\mathbb{D} = \mathbb{C}$ ($\dim_{\mathbb{R}} \mathbb{D} = 2$) for $p - q \equiv 3, 7 \pmod 8$;
\item $\mathbb{D} = \mathbb{H}$ ($\dim_{\mathbb{R}} \mathbb{D} = 4$) for $p - q \equiv 4, 5, 6 \pmod 8$.
\end{itemize}
In $M_N(\mathbb{D})$ ($N \ge 1$), an element $X$ lies in $Z(M_N(\mathbb{D})) = M_N(\mathbb{D}) \setminus \GL_N(\mathbb{D})$ iff its right $\mathbb{D}$-column rank is strictly less than $N$. For $N = 1$, $Z(\mathbb{D}) = \{0\}$ has real codimension $\dim_{\mathbb{R}} \mathbb{D} \in \{1, 2, 4\}$. For $N \ge 2$, the dense open stratum of corank-$1$ matrices near $\mathrm{diag}(I_{N-1}, 0)$ consists of block matrices $\bigl(\begin{smallmatrix} A & B \\ C & D \end{smallmatrix}\bigr)$ with $A \in \GL_{N-1}(\mathbb{D})$ satisfying the single Schur-complement equation $D - C A^{-1} B = 0$ in $\mathbb{D}$, which is a smooth real submanifold of real codimension $\dim_{\mathbb{R}} \mathbb{D}$, while higher-corank strata have codimension $k^2 \dim_{\mathbb{R}} \mathbb{D} > \dim_{\mathbb{R}} \mathbb{D}$ ($k \ge 2$). Passing to $M_N(\mathbb{D}) \oplus M_N(\mathbb{D})$ gives $Z = (Z(M_N(\mathbb{D})) \times M_N(\mathbb{D})) \cup (M_N(\mathbb{D}) \times Z(M_N(\mathbb{D})))$, preserving the real codimension $\dim_{\mathbb{R}} \mathbb{D} = c(p - q \bmod 8)$.
\end{proof}

\begin{theorem}[Exact dimension law over algebraically closed fields]\label{thm:complex-law}
Let $K$ be an algebraically closed field with $\mathrm{char}\,K \neq 2$ (for instance $K = \mathbb{C}$). For every integer $n \ge 0$ and every quadratic form $Q$ on $K^n$ of arbitrary metric rank deficit $r \in \{0, 1, \dots, n\}$,
\[
\dim_K Z\bigl(\Cl(K^n, Q)\bigr) = 2^n - 1, \qquad \codim_K Z\bigl(\Cl(K^n, Q)\bigr) = 1.
\]
\end{theorem}

\begin{proof}
For $N = 2^n \ge 1$, $P(x) := \det(L_x)$ is a homogeneous polynomial of degree $N$ on $\Cl(K^n, Q) \cong K^N$ satisfying $P(1) = 1 \neq 0$ and $P(0) = 0$. For $n = 0$, $Z(K) = \{0\}$ has dimension $0 = 2^0 - 1$. For $n \ge 1$, $P$ is nonconstant of positive degree over an algebraically closed field $K$, so by Krull's principal ideal theorem its zero locus $Z(\Cl(K^n, Q)) = V(P)$ has pure dimension $N - 1 = 2^n - 1$, regardless of $r$.
\end{proof}

\begin{corollary}[Refutation of Conjecture 00000007621]\label{cor:refutation}
Conjecture~00000007621 is false both over $\mathbb{R}$ and over $\mathbb{C}$:
\begin{enumerate}
\setlength{\itemsep}{2pt}\setlength{\parskip}{0pt}
\item[(a)] Over $\mathbb{R}$, no function $F(n, r)$ of the ambient dimension $n$ and the metric rank deficit $r$ can equal $\dim_{\mathbb{R}} Z(\Cl(V, Q))$, and the subtracted constant $c \in \{1, 2, 4\}$ in~\eqref{eq:bott-law} is $\dim_{\mathbb{R}} \mathbb{D}_{p-q \bmod 8}$ rather than $r$.
\item[(b)] Over $\mathbb{C}$, $\dim_{\mathbb{C}} Z(\Cl(\mathbb{C}^n, Q)) = 2^n - 1$ is strictly independent of $r \in \{0, \dots, n\}$, and its codimension $1$ disagrees with $r$ for every $r \neq 1$ (in particular for every nondegenerate metric $r = 0$).
\end{enumerate}
\end{corollary}

\section{Explicit Low-Dimensional and Degenerate Counterexamples}

Table~\ref{tab:clifford-dims} exhibits the complete zero-divisor classification across $n \in \{1, 2, 3\}$, making all obstructions concrete without appealing to Bott periodicity.

\begin{table}[ht]
\centering
\small
\begin{tabular}{@{}ccccclcc@{}}
\toprule
$n$ & $(p, q, r)$ & $\rank(B_Q)$ & Deficit $r$ & $\dim_{\mathbb{R}} A$ & Algebra $\Cl_{p,q,r}(\mathbb{R})$ & $\dim_{\mathbb{R}} Z(A)$ & $\codim_{\mathbb{R}} Z(A)$ \\
\midrule
$1$ & $(0, 1, 0)$ & $1$ & $0$ & $2$ & $\mathbb{C}$ & $0$ & $2$ \\
$1$ & $(1, 0, 0)$ & $1$ & $0$ & $2$ & $\mathbb{R} \oplus \mathbb{R}$ & $1$ & $1$ \\
$1$ & $(0, 0, 1)$ & $0$ & $1$ & $2$ & $\Lambda(\mathbb{R}^1) \cong \mathbb{R}[\varepsilon]/(\varepsilon^2)$ & $1$ & $1$ \\
\midrule
$2$ & $(0, 2, 0)$ & $2$ & $0$ & $4$ & $\mathbb{H}$ & $0$ & $4$ \\
$2$ & $(1, 1, 0)$ & $2$ & $0$ & $4$ & $M_2(\mathbb{R})$ & $3$ & $1$ \\
$2$ & $(2, 0, 0)$ & $2$ & $0$ & $4$ & $M_2(\mathbb{R})$ & $3$ & $1$ \\
$2$ & $(0, 1, 1)$ & $1$ & $1$ & $4$ & $\mathbb{C} \ltimes \mathbb{C} g_1$ ($g_1^2 = 0$) & $2$ & $2$ \\
$2$ & $(1, 0, 1)$ & $1$ & $1$ & $4$ & $(\mathbb{R} \oplus \mathbb{R}) \ltimes (\mathbb{R} \oplus \mathbb{R}) g_1$ & $3$ & $1$ \\
$2$ & $(0, 0, 2)$ & $0$ & $2$ & $4$ & $\Lambda(\mathbb{R}^2)$ & $3$ & $1$ \\
\midrule
$3$ & $(0, 3, 0)$ & $3$ & $0$ & $8$ & $\mathbb{H} \oplus \mathbb{H}$ & $4$ & $4$ \\
$3$ & $(1, 2, 0)$ & $3$ & $0$ & $8$ & $M_2(\mathbb{C})$ & $6$ & $2$ \\
$3$ & $(3, 0, 0)$ & $3$ & $0$ & $8$ & $M_2(\mathbb{C})$ & $6$ & $2$ \\
$3$ & $(2, 1, 0)$ & $3$ & $0$ & $8$ & $M_2(\mathbb{R}) \oplus M_2(\mathbb{R})$ & $7$ & $1$ \\
\bottomrule
\end{tabular}
\caption{Exact dimensions of $Z(\Cl_{p,q,r}(\mathbb{R}))$ for $n \le 3$. Even among nondegenerate metrics ($r = 0$), $\dim_{\mathbb{R}} Z$ takes two values at $n = 1$, two values at $n = 2$, and three values $\{4, 6, 7\}$ at $n = 3$.}
\label{tab:clifford-dims}
\end{table}

\begin{proposition}[Direct verification for $n = 1, 2$ and $\Lambda(K^r)$]\label{prop:direct}
Let $K$ be a field.
\begin{enumerate}
\setlength{\itemsep}{2pt}\setlength{\parskip}{0pt}
\item[(i)] \textbf{Division vs.\ split at $(n, r) = (1, 0)$:} In $\Cl_{0,1}(K) = K \oplus K e_1$ ($e_1^2 = -1$), $(a + b e_1)(a - b e_1) = a^2 + b^2$; if $a^2 + b^2 = 0 \implies a = b = 0$ (e.g.\ $K = \mathbb{R}, \mathbb{Q}$, or $\mathbb{F}_q$ with $q \equiv 3 \pmod 4$), then $Z(\Cl_{0,1}(K)) = \{0\}$ ($\dim_K Z = 0$). In $\Cl_{1,0}(K) = K \oplus K e_1$ ($e_1^2 = +1$), $a + b e_1 \in Z(\Cl_{1,0}(K)) \iff a^2 - b^2 = 0 \iff a = \pm b$, annihilated by $1 \mp e_1 \neq 0$, so $Z(\Cl_{1,0}(K)) = K(1 + e_1) \cup K(1 - e_1)$ ($\dim_K Z = 1$).
\item[(ii)] \textbf{Quaternions vs.\ matrices at $(n, r) = (2, 0)$:} In $\Cl_{0,2}(K) \cong \mathbb{H}_K$, $q \bar{q} = \bar{q} q = a^2 + b^2 + c^2 + d^2$; if $K$ is $4$-square anisotropic (e.g.\ $K = \mathbb{R}, \mathbb{Q}$), then $Z(\Cl_{0,2}(K)) = \{0\}$ ($\dim_K Z = 0$). In $\Cl_{1,1}(K)$ ($e_1^2 = 1, e_2^2 = -1$), every $q = a + b e_1 + c e_2 + d e_1 e_2$ satisfies $q \tilde{q} = \tilde{q} q = (a^2 - b^2 + c^2 - d^2)\cdot 1$ where $\tilde{q} = a - b e_1 - c e_2 - d e_1 e_2$, so $q \in Z(\Cl_{1,1}(K)) \iff a^2 - b^2 + c^2 - d^2 = 0$ (matching $Z(M_2(K)) = \{X \in M_2(K) : \det X = 0\}$, $\dim_K Z = 3$).
\item[(iii)] \textbf{Purely degenerate exterior algebra $\Lambda(K^r) = \Cl_{0,0,r}(K)$ ($n = r \ge 1$):} Writing $x = x_0 + x_+ \in K \oplus \Lambda^+(K^r)$, one has $x \in Z(\Lambda(K^r)) \iff x_0 = 0$: if $x_0 = 0$, then $x \wedge (e_1 \wedge \dots \wedge e_r) = 0$ with $e_1 \wedge \dots \wedge e_r \neq 0$, whereas if $x_0 \neq 0$, then $x = x_0(1 + x_0^{-1} x_+)$ has inverse $x_0^{-1}\sum_{m=0}^r (-x_0^{-1} x_+)^m$. Hence $\dim_K Z(\Lambda(K^r)) = 2^r - 1$ ($\codim_K Z = 1$), identical to split nondegenerate $\Cl_{m,m}(K)$ ($r = 0$).
\end{enumerate}
\end{proposition}

\section{Alternative Readings: Vectors, Annihilators, Finite Fields, and Pairs}

To verify that no alternative algebraic or geometric reading rescues the conjecture, we analyze four further natural formulations.

\begin{proposition}[Grade-$1$ vector zero divisors and the isotropic cone]\label{prop:grade1}
For any quadratic space $(V, Q)$ over a field $K$, a grade-$1$ vector $v \in V \subset \Cl(V, Q)$ lies in $Z(\Cl(V, Q))$ if and only if $Q(v) = 0$:
\[
Z\bigl(\Cl(V, Q)\bigr) \cap V = \{v \in V : Q(v) = 0\}.
\]
In particular, at $(n, r) = (2, 0)$ over $\mathbb{R}$, $\dim_{\mathbb{R}}(Z(\Cl_{0,2}(\mathbb{R})) \cap \mathbb{R}^2) = \dim_{\mathbb{R}}(Z(\Cl_{2,0}(\mathbb{R})) \cap \mathbb{R}^2) = 0$, whereas $\dim_{\mathbb{R}}(Z(\Cl_{1,1}(\mathbb{R})) \cap \mathbb{R}^2) = \dim_{\mathbb{R}}\{(x, y) \in \mathbb{R}^2 : x^2 - y^2 = 0\} = 1 \neq 0 = r$.
\end{proposition}

\begin{proof}
In $\Cl(V, Q)$, $v^2 = Q(v)\cdot 1$. If $Q(v) \neq 0$, then $v \cdot (Q(v)^{-1} v) = (Q(v)^{-1} v) \cdot v = 1$, so $v \in \Cl(V, Q)^\times$ and $v \notin Z(\Cl(V, Q))$. Conversely, if $Q(v) = 0$, then either $v = 0 \in Z$ or $v \neq 0$ annihilates itself ($v \cdot v = 0$), so $v \in Z(\Cl(V, Q)) \cap V$.
\end{proof}

\begin{remark}[Diagnostic origin of the misconception]
Proposition~\ref{prop:grade1} pinpoints a plausible heuristic behind the false conjecture: if one restricts simultaneously to \emph{grade-$1$ vectors} $V \subset \Cl(V, Q)$ and to \emph{semidefinite real metrics} ($q = 0$ or $p = 0$), then $\{v \in \mathbb{R}^n : Q(v) = 0\} = \ker(B_Q)$ indeed has dimension $r$. That equality fails immediately upon allowing mixed-grade algebra elements ($1 \pm e_1 \in Z(\Cl_{1,0}(\mathbb{R}))$ at $q = r = 0$), indefinite real metrics ($x^2 - y^2 = 0$ at $r = 0$), or complex scalars.
\end{remark}

\begin{proposition}[Multiplicative annihilator nullity $\dim_K \Ann_R(x)$]\label{prop:annihilator}
Let $K$ be any field.
\begin{enumerate}
\setlength{\itemsep}{2pt}\setlength{\parskip}{0pt}
\item[(i)] In $M_2(K)$ ($n = 2$, metric nullity $r = 0$ when $\mathrm{char}\,K \neq 2$), the nonzero zero divisor $E_{11} = \bigl(\begin{smallmatrix} 1 & 0 \\ 0 & 0 \end{smallmatrix}\bigr)$ has $\Ann_R(E_{11}) = \bigl\{\bigl(\begin{smallmatrix} 0 & 0 \\ z & w \end{smallmatrix}\bigr) : z, w \in K\bigr\}$ of dimension $2 \neq 0 = r$.
\item[(ii)] In $\Lambda(K^r)$ for every $r \ge 2$, $\dim_K \Ann_R(x)$ is non-constant across $x \in Z(\Lambda(K^r)) \setminus \{0\}$:
\[
\dim_K \Ann_R(e_r) = 2^{r-1} \;<\; 2^r - 1 = \dim_K \Ann_R(e_1 \wedge \dots \wedge e_r).
\]
\end{enumerate}
\end{proposition}

\begin{proof}
Part~(i) is immediate from $\bigl(\begin{smallmatrix} 1 & 0 \\ 0 & 0 \end{smallmatrix}\bigr)\bigl(\begin{smallmatrix} x & y \\ z & w \end{smallmatrix}\bigr) = \bigl(\begin{smallmatrix} x & y \\ 0 & 0 \end{smallmatrix}\bigr)$. For~(ii), decompose $\Lambda(K^r) = \Lambda(K^{r-1}) \oplus \Lambda(K^{r-1}) \wedge e_r$. Writing $y = u + v \wedge e_r$, we have $e_r \wedge y = \alpha(u) \wedge e_r = 0 \iff u = 0$, so $\Ann_R(e_r) = \Lambda(K^{r-1}) \wedge e_r$ has dimension $2^{r-1}$. Meanwhile $(e_1 \wedge \dots \wedge e_r) \wedge y = y_0 (e_1 \wedge \dots \wedge e_r) = 0 \iff y_0 = 0$, so $\Ann_R(e_1 \wedge \dots \wedge e_r) = \Lambda^+(K^r)$ has dimension $2^r - 1$. For every $r \ge 2$, $2^{r-1} < 2^r - 1$.
\end{proof}

\begin{proposition}[Exact finite-field point-count polynomials over $\mathbb{F}_q$]\label{prop:fq}
Let $\mathbb{F}_q$ be a finite field of cardinality $q$.
\begin{enumerate}
\setlength{\itemsep}{2pt}\setlength{\parskip}{0pt}
\item[(i)] If $q \equiv 3 \pmod 4$, then $-1 \notin (\mathbb{F}_q^\times)^2$, so $\Cl_{0,1}(\mathbb{F}_q) \cong \mathbb{F}_{q^2}$ has $|Z(\Cl_{0,1}(\mathbb{F}_q))| = 1$ (degree $0$ in $q$), whereas $|Z(\Cl_{1,0}(\mathbb{F}_q))| = 2q - 1$ (degree $1$ in $q$), both at $n = 1, r = 0$.
\item[(ii)] Over every $\mathbb{F}_q$, $|Z(M_2(\mathbb{F}_q))| = q^4 - |\GL_2(\mathbb{F}_q)| = q^4 - (q^2 - 1)(q^2 - q) = q^3 + q^2 - q$ (degree $3$ in $q$ at $r = 0$).
\item[(iii)] For every $r \ge 0$, $|Z(\Lambda(\mathbb{F}_q^r))| = q^{2^r - 1}$ (degree $2^r - 1$ in $q$); at $r = 2$, $|Z(\Lambda(\mathbb{F}_q^2))| = q^3$ has the same degree $3$ as $|Z(M_2(\mathbb{F}_q))|$ ($r = 0$).
\end{enumerate}
\end{proposition}

\begin{proof}
Claim~(i) follows from Proposition~\ref{prop:direct}(i): $(0, 0)$ is unique when $-1$ is a non-square, whereas $a = \pm b$ with $\mathrm{char}\,\mathbb{F}_q \neq 2$ yields $1 + 2(q - 1) = 2q - 1$ solutions. For~(ii), partitioning $\bigl(\begin{smallmatrix} a & b \\ c & d \end{smallmatrix}\bigr)$ with $ad - bc = 0$ into $(a, b) = (0, 0)$ ($q^2$ matrices), $a \neq 0$ ($(q - 1)q^2$ matrices with $d = a^{-1}bc$), and $a = 0, b \neq 0$ ($(q - 1)q$ matrices with $c = 0$) gives $q^2 + (q - 1)q^2 + (q - 1)q = q^3 + q^2 - q = q^4 - (q^2 - 1)(q^2 - q)$. Claim~(iii) follows from $Z(\Lambda(\mathbb{F}_q^r)) = \{x : x_0 = 0\} \cong \mathbb{F}_q^{2^r - 1}$.
\end{proof}

\begin{remark}[Incidence pairs and commutative Jordan spin factors]
Even if one replaces $Z(A)$ by the nontrivial incidence variety $\mathcal{Z}_2^\circ(A) = \{(x, y) \in (A \setminus \{0\})^2 : xy = 0\}$, at $n = 2, r = 0$ one has $\mathcal{Z}_2^\circ(\mathbb{H}) = \emptyset$ while $\dim_{\mathbb{R}} \overline{\mathcal{Z}_2^\circ(M_2(\mathbb{R}))} = 3 + 2 = 5 = \dim_{\mathbb{R}} \overline{\mathcal{Z}_2^\circ(\Lambda(\mathbb{R}^2))}$ ($r = 2$). Likewise, in the commutative Jordan spin factor $J(V, Q) = K\cdot 1 \oplus V$ with product $(a + u) \circ (b + v) = (ab + B_Q(u, v)) + (av + bu)$, already at $n = 1, r = 0$ one has $J(\mathbb{R}, \pm u^2) = \Cl_{1,0}(\mathbb{R})$ vs.\ $\Cl_{0,1}(\mathbb{R})$ ($\dim_{\mathbb{R}} Z = 1$ vs.\ $0$).
\end{remark}

\section{Formal Verification in Core Lean 4.19.0 and Computational Checks}

The disproof is accompanied by a self-contained Core Lean~4.19.0 library (\texttt{lean/}) built with \texttt{-DwarningAsError=true}, zero external dependencies, zero \texttt{sorry}, zero \texttt{native\_decide}, and only foundational Lean axioms (\texttt{propext}, \texttt{Classical.choice}, \texttt{Quot.sound}). Table~\ref{tab:lean-map} summarizes the principal machine-checked theorems. An independent Python~3.12 verification suite (\texttt{verify.py}) additionally cross-checks associativity, exact Gaussian-elimination ranks over $\mathbb{Q}$ and $\mathbb{F}_p$, radical projection invariance, isotropic cones, annihilator nullities, and $\mathbb{F}_p$ point counts across all 35 signature triples with $n \le 4$.

\begin{table}[ht]
\centering
\footnotesize
\begin{tabular}{@{}ll@{}}
\toprule
Mathematical Result & Core Lean 4.19.0 Declaration \\
\midrule
Orthogonal Clifford step ($\varepsilon \in \{0, \pm 1\}$) & \texttt{AssocInvolAlg.cliffordStep} \\
Master radical reduction ($x \in Z \leftrightarrow \pi(x) \in Z$) & \texttt{AssocInvolAlg.radIter\_isZeroDivisor\_iff} \\
Exterior algebra $Z(\Lambda(K^r)) = \{x_0 = 0\}$ & \texttt{AssocInvolAlg.extAlg\_isZeroDivisor\_iff} \\
Annihilator non-constancy in $\Lambda(K^r)$ ($r \ge 2$) & \texttt{AssocInvolAlg.extAlg\_annihilator\_nonconstant} \\
Division vs.\ split at $(n, r) = (1, 0)$ & \texttt{cl01\_isZeroDivisor\_iff}, \texttt{cl10\_isZeroDivisor\_iff} \\
Quaternions vs.\ split quadric at $(n, r) = (2, 0)$ & \texttt{cl02\_isZeroDivisor\_iff}, \texttt{cl11\_isZeroDivisor\_iff} \\
Grade-$1$ isotropic cone in $\Cl_{1,1}(K)$ ($r = 0$) & \texttt{cl11\_grade1\_isZeroDivisor\_iff} \\
Grade-$1$ anisotropic cone in $\Cl_{0,2}(K)$ ($r = 0$) & \texttt{cl02\_grade1\_isZeroDivisor\_iff} \\
Matrix determinant law $Z(M_2(K)) = \{\det = 0\}$ & \texttt{Mat2.isZeroDivisor\_iff\_det\_eq\_zero} \\
Nonzero annihilator $\Ann_R(E_{11})$ at $r = 0$ & \texttt{Mat2.inRightAnn\_e11\_iff} \\
Exact $\mathbb{F}_q$ counts ($1$, $2q-1$, $q^3+q^2-q$, $q^{2^r-1}$) & \texttt{FiniteFieldOps.mat2\_zeroDivisor\_card} \\
Exterior \& annihilator $\mathbb{F}_q$ counts ($q^{2^r-1}, q^{2^{r-1}}$) & \texttt{FiniteFieldOps.extAlg\_zeroDivisor\_card} \\
End-to-end disproof over $\mathbb{Q}$ and $\mathbb{F}_3$ & \texttt{conjecture\_00000007621\_disproved} \\
\bottomrule
\end{tabular}
\caption{Correspondence between mathematical theorems in Sections~1--4 and formal declarations in \texttt{ZeroDivisorMetricRankDeficit}.}
\label{tab:lean-map}
\end{table}

\end{document}
